% ===========================================================================
%  Stil des Verteidigungsvortrags -- von vortrag.tex geladen.
%
%  Leitgedanke: EINE Aussage pro Folie, gross gesetzt.  Die Grundschrift ist
%  deshalb die 14pt-Klasse und nicht 10 oder 11pt; Aufzaehlungen laufen in
%  \large.  Was hier nicht als Makro vorgesehen ist, gehoert nicht auf eine
%  Folie, sondern in den Sprechtext (\note) oder auf eine Zusatzfolie.
%
%  Die mitgelieferten beamer-Themen bringen Navigationsleisten und
%  Schattenkaesten mit, die auf einem Projektor nur Platz kosten -- deshalb
%  das Standardthema als Grundlage und alles Sichtbare hier definiert.
% ===========================================================================
\usetheme{default}
\usecolortheme{default}
\setbeamertemplate{navigation symbols}{}

\usepackage{iftex}
\ifPDFTeX
  \usepackage[utf8]{inputenc}
  \usepackage[T1]{fontenc}
  \usepackage{lmodern}
\else
  \usepackage{fontspec}
\fi
\usepackage[ngerman]{babel}
\usepackage[german=quotes]{csquotes}
\usepackage{amsmath,amssymb}
\usepackage{booktabs}
\usepackage{array}
\usepackage{graphicx}
\usepackage{tikz}
\usetikzlibrary{arrows.meta,positioning,calc,fit}
\usepackage{appendixnumberbeamer}

% --------------------------------------------------------------------------
%  Farben.  akzent traegt Titel und Hervorhebung, gut/warn die Richtung eines
%  Effekts, die Belegart-Marker die Beweislage.
% --------------------------------------------------------------------------
\definecolor{akzent}{RGB}{20,64,116}
\definecolor{akzenthell}{RGB}{72,124,180}
\definecolor{gut}{RGB}{22,110,60}
\definecolor{warn}{RGB}{176,58,20}
\definecolor{grau}{RGB}{115,115,115}
\definecolor{hgrau}{RGB}{243,244,246}

\setbeamercolor{structure}{fg=akzent}
\setbeamercolor{frametitle}{fg=akzent,bg=}
\setbeamercolor{framesubtitle}{fg=grau,bg=}
\setbeamercolor{footline}{fg=grau}
\setbeamercolor{alerted text}{fg=warn}

\setbeamerfont{frametitle}{size=\Large,series=\bfseries}
\setbeamerfont{framesubtitle}{size=\normalsize,series=\mdseries}
\setbeamerfont{footline}{size=\tiny}
\setbeamerfont{itemize/enumerate body}{size=\large}
\setbeamerfont{itemize/enumerate subbody}{size=\normalsize}
\setbeamerfont{description item}{series=\bfseries}

% Flache Striche statt Kugeln: laeuft bei grosser Schrift ruhiger.
\newcommand{\markerI}{\color{akzenthell}\rule[0.30ex]{0.55em}{1.1pt}}
\newcommand{\markerII}{\color{grau}\rule[0.30ex]{0.38em}{0.8pt}}
\setbeamertemplate{itemize item}{\markerI}
\setbeamertemplate{itemize subitem}{\markerII}
\setbeamertemplate{enumerate item}{\color{akzent}\bfseries\insertenumlabel.}

\setbeamersize{text margin left=1.1em,text margin right=1.1em}

% Die Aufzaehlung des Hauptteils.  Bei 14pt-Grundschrift kleben Punkte mit
% dem beamer-Standardabstand aneinander; enumitem waere der uebliche Weg,
% bricht aber beamers \insertenumlabel (\labelenumi -> {\labelenumi},
% Endlosrekursion).  Deshalb eine eigene Liste: der Abstand ist damit exakt
% bekannt, und genau darauf stuetzt sich das Hoehenbudget weiter unten.
%   \begin{punkte}          Hauptteil, luftig
%   \begin{punkte}[0.25em]  Zusatzfolien, dicht
\newenvironment{punkte}[1][0.55em]{%
  \begin{list}{\markerI}{%
    \setlength{\leftmargin}{1.6em}%
    \setlength{\labelwidth}{0.8em}%
    \setlength{\labelsep}{0.7em}%
    \setlength{\itemsep}{#1}%
    \setlength{\parsep}{0pt}%
    \setlength{\topsep}{0.4em}%
    \setlength{\partopsep}{0pt}%
  }}{\end{list}}

% --------------------------------------------------------------------------
%  Folientitel mit Trennlinie, Fusszeile mit Abschnitt und Foliennummer.
% --------------------------------------------------------------------------
\makeatletter
\setbeamertemplate{frametitle}{%
  \nointerlineskip\vskip0.85em%
  \begin{beamercolorbox}[wd=\paperwidth,leftskip=\beamer@leftmargin,
                         rightskip=\beamer@rightmargin]{frametitle}%
    \usebeamerfont{frametitle}\insertframetitle%
    \ifx\insertframesubtitle\@empty\else%
      \\[0.2em]{\usebeamerfont{framesubtitle}%
        \usebeamercolor[fg]{framesubtitle}\insertframesubtitle}%
    \fi%
    \\[0.3em]{\color{akzent!22}\rule{\dimexpr\paperwidth-\beamer@leftmargin-\beamer@rightmargin\relax}{1pt}}%
  \end{beamercolorbox}%
  \vskip0.6em%
}
\setbeamertemplate{footline}{%
  \begin{beamercolorbox}[wd=\paperwidth,ht=2.4ex,dp=1.6ex,
                         leftskip=\beamer@leftmargin,
                         rightskip=\beamer@rightmargin]{footline}%
    \usebeamerfont{footline}%
    \insertshorttitle\hfill\insertsectionhead\hfill%
    \insertframenumber\,/\,\inserttotalframenumber%
  \end{beamercolorbox}%
}
\makeatother

% --------------------------------------------------------------------------
%  Makros -- dieselben Namen wie in praeambel.tex, damit Textbausteine und
%  Zahlen aus kapitel/ unveraendert uebernommen werden koennen.
% --------------------------------------------------------------------------
\newcommand{\code}[1]{\texttt{#1}}
\newcommand{\px}{\,px}
\newcommand{\PP}{\,\textnormal{\%-Punkte}}  % \textnormal, damit der Bindestrich
                                            % im Mathemodus kein Minus wird
\newcommand{\mAPfifty}{mAP@0{,}5}
\newcommand{\mAPseven}{mAP@0{,}75}
\newcommand{\mAPall}{mAP@0{,}5{:}0{,}95}
\newcommand{\Oe}{\ensuremath{\varnothing}}
\newcolumntype{L}[1]{>{\raggedright\arraybackslash}p{#1}}
\newcolumntype{R}[1]{>{\raggedleft\arraybackslash}p{#1}}

% Belegart: M gemessen, P physikalisch, L Literatur, A Annahme.
\newcommand{\bM}{\textcolor{gut}{\textbf{M}}}
\newcommand{\bP}{\textcolor{akzent}{\textbf{P}}}
\newcommand{\bL}{\textcolor{akzenthell}{\textbf{L}}}
\newcommand{\bA}{\textcolor{warn}{\textbf{A}}}

% --------------------------------------------------------------------------
%  Das Hoehenbudget einer Folie -- an dieser Klasse gemessen, nicht geschaetzt.
%
%  \textheight ist 241 pt, \textwidth 425 pt (14,97 cm).  Davon gehen ab:
%    Folientitel ohne Untertitel   ~50 pt      Fusszeile   ~14 pt
%    Folientitel mit  Untertitel   ~71 pt
%  Es bleiben also 177 pt (ohne Untertitel) bzw. 156 pt (mit Untertitel).
%
%  Was das kostet:
%    ein Aufzaehlungspunkt (\large, einzeilig)   22 pt  + 9 pt Abstand
%    eine Umbruchzeile darin                     22 pt
%    \fuss{}                                    ~34 pt
%    \grosszahl{}{}{}                           ~124 pt
%    \zweizahl{}{}{}{}                           ~89 pt
%    \botschaft{} je Zeile                       ~39 pt
%
%  Daraus die Obergrenzen, die beim Schreiben gelten:
%    mit Untertitel     4 Punkte   ODER 3 Punkte + \fuss
%    ohne Untertitel    5 Punkte   ODER 4 Punkte + \fuss
%    Grosszahl + \fuss  nur ohne Untertitel
%
%  Wer mehr braucht, hat zwei Folien -- nicht eine kleinere Schrift.
%  'make vortrag-pruefen' nennt jede Folie, die die Grenze reisst.
% --------------------------------------------------------------------------

% --------------------------------------------------------------------------
%  Folientypen.  Jeder steht fuer genau eine Art von Aussage; zusammen sind
%  sie die einzigen erlaubten Bauformen einer Folie im Hauptteil.
% --------------------------------------------------------------------------

% (1) Botschaft: ein Satz, der die Folie ganz fuellt.  Fuer die Kernsaetze.
\newcommand{\botschaft}[1]{%
  \vfill
  {\LARGE\baselineskip=1.35\baselineskip\color{akzent}#1\par}
  \vfill}

% (2) Grosszahl: die Zahl ist der Inhalt, nicht ihre Tabellenzeile.
%     #1 Wert, #2 Beschriftung, #3 Zusatz (darf leer sein)
\newcommand{\grosszahl}[3]{%
  \begin{center}
    {\fontsize{62}{66}\selectfont\bfseries\color{akzent}#1}\\[0.35em]
    {\Large #2}%
    \ifx\relax#3\relax\else\\[0.2em]{\normalsize\color{grau}#3}\fi
  \end{center}}

% (3) Zwei Zahlen nebeneinander -- fuer val gegen test.
\newcommand{\zweizahl}[4]{%
  \begin{center}
  \begin{minipage}{0.46\textwidth}\centering
    {\fontsize{56}{62}\selectfont\bfseries\color{akzent}#1}\\[0.35em]{\large #2}
  \end{minipage}\hfill
  \begin{minipage}{0.46\textwidth}\centering
    {\fontsize{56}{62}\selectfont\bfseries\color{akzent}#3}\\[0.35em]{\large #4}
  \end{minipage}
  \end{center}}

% (3b) Drei Zahlen -- fuer das verfuegbare Quellmaterial 9 / 7 / 2.
\newcommand{\dreizahl}[6]{%
  \begin{center}
  \begin{minipage}{0.30\textwidth}\centering
    {\fontsize{54}{60}\selectfont\bfseries\color{akzent}#1}\\[0.3em]{\large #2}
  \end{minipage}\hfill
  \begin{minipage}{0.30\textwidth}\centering
    {\fontsize{54}{60}\selectfont\bfseries\color{akzent}#3}\\[0.3em]{\large #4}
  \end{minipage}\hfill
  \begin{minipage}{0.30\textwidth}\centering
    {\fontsize{54}{60}\selectfont\bfseries\color{akzent}#5}\\[0.3em]{\large #6}
  \end{minipage}
  \end{center}}

% (3c) Von-nach-Zahl: der Effekt ist der Sprung, nicht der Endwert.
\newcommand{\sprungzahl}[5]{%
  \begin{center}
    {\Large\color{grau}#1}\quad
    {\fontsize{42}{46}\selectfont\bfseries\color{grau}#2}\;
    {\fontsize{42}{46}\selectfont\color{akzenthell}$\rightarrow$}\;
    {\fontsize{42}{46}\selectfont\bfseries\color{akzent}#3}\\[0.4em]
    {\Large #4}%
    \ifx\relax#5\relax\else\\[0.2em]{\normalsize\color{grau}#5}\fi
  \end{center}}

% (3d) Schlagwortreihe: sechs Begriffe als Inhalt, ohne Erklaerung.
%      Die Erklaerung steht im Sprechtext und auf der Zusatzfolie.
\newcommand{\schlagworte}[1]{%
  \begin{center}\large
    \begingroup\setlength{\tabcolsep}{0.55em}#1\endgroup
  \end{center}}

% (4) Fussnotiz einer Folie: die eine Einschraenkung, die dazugehoert.
%     Bewusst klein und grau -- sie ist Beiwerk, nicht Inhalt.
% Regel: hoechstens zwei Zeilen.  Wer drei braucht, gehoert in den
% Sprechtext -- \fuss traegt eine Einschraenkung, nie ein Argument.
\newcommand{\fuss}[1]{%
  \par\vspace{0.5em}
  {\color{akzent!18}\rule{0.30\linewidth}{0.7pt}}\par\vspace{0.25em}
  {\footnotesize\color{grau}#1\par}}

% (4b) Abbildung mit Hoehendeckel.  \bild deckelt auf eine Hoehe, die mit
%      Titel UND \fuss noch passt; \bildgross auf eine, die ohne \fuss passt.
%      keepaspectratio, damit breite Plots an der Breite und hohe Fotos an
%      der Hoehe begrenzt werden.
\newcommand{\bild}[1]{%
  \begingroup\centering
    \includegraphics[width=\textwidth,height=0.50\textheight,
                     keepaspectratio]{#1}\par
  \endgroup}
\newcommand{\bildgross}[1]{%
  \begingroup\centering
    \includegraphics[width=\textwidth,height=0.54\textheight,
                     keepaspectratio]{#1}\par
  \endgroup}

% (5) Abschnittstrenner.
\newcommand{\trenner}[2]{%
  \begin{frame}[plain,noframenumbering]
    \vfill
    \begin{center}
      {\color{grau}\large #1}\\[0.5em]
      {\color{akzent}\Huge\bfseries #2}
    \end{center}
    \vfill
  \end{frame}}

% (6) Hervorgehobener Kasten fuer eine einzelne kurze Aussage.
\newcommand{\kasten}[1]{%
  \begingroup\setlength{\fboxsep}{0.7em}%
  \noindent\colorbox{hgrau}{\begin{minipage}{\dimexpr\linewidth-1.4em\relax}%
    \centering#1%
  \end{minipage}}\endgroup}

% (7) Budgetmarker.  Alle drei Entwuerfe nannten dasselbe Risiko: Auf den
%     Ablationsfolien liegen 0,961 und 0,969 nur wenige Folien auseinander,
%     stammen aber aus verschiedenen Budgets.  Der Marker steht deshalb auf
%     JEDER Ablationsfolie an derselben Stelle.
\newcommand{\budgetmarker}{%
  \par\nointerlineskip
  \hfill\raisebox{0pt}[0pt][0pt]{%
    \colorbox{hgrau}{\footnotesize\color{grau}%
      \,3000 Bilder $\cdot$ 30 Epochen $\cdot$ nicht mit dem Referenzmodell
      vergleichbar\,}}\par\nointerlineskip}

% Zusatzfolien im Anhang duerfen dicht sein -- sie werden nur in der
% Diskussion gezeigt und dann gelesen, nicht vorgetragen.
% Referenzmaterial, aus der Naehe gelesen -- nicht projizierte Prosa.
\newcommand{\zusatzdicht}{\fontsize{7.5}{9}\selectfont}
